% SPDX-License-Identifier: Apache-2.0
\section{How Chrome Remote Desktop forwards WebAuthn}
\label{app:crd}

Chrome Remote Desktop (CRD) has its own way to use a security key from a
remote host. It does not forward the key. It forwards each WebAuthn request
to the client, and the client's own Chrome runs it. On a Chromebook the key
then stays with ChromeOS: no Linux container, no root, and no setuid program
on the host. Google opens this path to \term{managed} accounts only. This
appendix explains how it works, what an organization must configure, and
what we could not verify. The Chromium sources are at commit
\code{2b88984cf66d} (2026-11-10)~\cite{crdhost,crdclient}.

\begin{figure*}[t]
\centering
\begin{tikzpicture}[font=\scriptsize]
  \tikzset{n/.style={box,text width=24mm,minimum height=9mm}}
  \node[sbox,n] (pg) at (0,0) {Web page in\\Google Chrome};
  \node[act,n] (ext) at (4.6,0) {CRD Security Key\\extension};
  \node[act,n] (nm) at (9.2,0) {\code{remote-webauthn}\\native host};
  \node[act,n] (h) at (13.8,0) {CRD host\\process};
  \draw[fl] (pg) -- node[above]{proxy} node[below]{API} (ext);
  \draw[fl] (ext) -- node[above]{native} node[below]{messaging} (nm);
  \draw[fl] (nm) -- node[above]{Mojo IPC} (h);
  \node[act,n] (cl) at (13.8,-2.6) {CRD web client\\(a web page)};
  \node[sbox,n] (wa) at (7.7,-2.6) {ChromeOS Chrome\\WebAuthn};
  \node[mem,n] (yk) at (2.3,-2.6) {YubiKey};
  \draw[fl] (h) -- node[right,align=left]{WebRTC data channel\\\code{remote-webauthn}} (cl);
  \draw[fl] (cl) -- node[above]{WebAuthn call with} node[below]{\code{remoteClientDataJSON}} (wa);
  \draw[fl] (wa) -- node[above]{USB HID,} node[below]{native} (yk);
  \begin{scope}[on background layer]
    \node[draw,dashed,rounded corners,inner sep=6pt,fit=(pg)(h),
          label={[font=\scriptsize]above:Remote host (Linux, CRD session)}] {};
    \node[draw,dashed,rounded corners,inner sep=6pt,fit=(yk)(cl),
          label={[font=\scriptsize]below:Chromebook (enrolled; user and device in the same organization)}] {};
  \end{scope}
\end{tikzpicture}
\caption{Remote WebAuthn in CRD. A request in Chrome on the host goes to the
client, and the client's Chrome runs it against the key with the site's
original origin. The response goes back the same way. Compare with
Fig.~\ref{fig:flow}, where the key itself is forwarded.}
\label{fig:crd}
\end{figure*}

\subsection{The host side}

Chrome on the host must have the \term{Chrome Remote Desktop Security Key
Extension}~\cite{crdext}. It uses the \code{webAuthenticationProxy}
extension API~\cite{wap}. When a CRD session starts, the CRD host writes a
file into a directory that Chrome watches. The extension then calls
\code{attach()}, and from that time Chrome sends every WebAuthn call to the
extension, not to local keys~\cite{crdhost} (verified). The extension passes
each request through native messaging to the program
\code{remote-webauthn}. That program passes it through Mojo IPC to the CRD
host process. The host sends it to the client over a WebRTC data channel
named \code{remote-webauthn}. There are four operations: create, get,
cancel, and a check for a platform authenticator. The request bodies are the
WebAuthn options as JSON~\cite{crdhost} (verified).

The CRD host turns this on for every host, by default; no policy is
needed~\cite{crdhost} (verified). Two conditions apply. On Linux,
\code{remote-webauthn} accepts only a caller that is Google Chrome from
\code{/opt/google/chrome}, not Chromium. And the CRD host does not install
the extension: an administrator must force-install it by policy, or the
user installs it~\cite{crdhost} (verified).

\subsection{The client side}

The CRD client is a web page. It calls WebAuthn in its own Chrome, and adds
the site's origin from the host in the extension \code{remoteClientDataJSON}
(earlier: \code{remoteDesktopClientOverride})~\cite{webauthn-ed} (verified).
Chrome then runs the ceremony as if the remote site had called it, with the
remote site's RP ID. The key signs client data that names the remote site,
so the site accepts the response.

This extension lets one page act for any site, so Chrome allows it only in
two cases~\cite{crdclient} (verified):
\begin{itemize}
\item The profile is \term{affiliated}, and the page's origin is in the
      policy \code{WebAuthentication}\allowbreak\code{RemoteDesktop}\allowbreak\code{AllowedOrigins}
      (Chrome~136 and later)~\cite{crdallow}. Affiliated means that the
      same organization manages the user and the device or browser.
\item A Google-internal policy is set, and the page is Google's internal
      CRD client. The source says that this case ``explicitly does not
      cover external instances of CRD''~\cite{crdclient}.
\end{itemize}
A personal \code{@gmail.com} account is not affiliated, so the first case
cannot apply to it. That is why this study built \code{ffwd-usbfs}.

\subsection{What a managed organization must configure}

Table~\ref{tab:crd} lists the settings for a Linux host and a Chromebook
client. One public report says that it works on ChromeOS from Chrome~147,
and that the Admin console has the setting since about
2026-07-14~\cite{iwadev}; we did not verify this against Google's own
documentation.

\begin{table}[t]
\centering
\caption{Settings for remote WebAuthn in CRD}
\label{tab:crd}
\scriptsize
\begin{tabular}{@{}>{\raggedright\arraybackslash}p{9mm}>{\raggedright\arraybackslash}p{52mm}>{\raggedright\arraybackslash}p{14mm}@{}}
\toprule
Where & Setting & Status \\
\midrule
Client & Chromebook enrolled in the organization that manages the user, so the profile is affiliated & Verified \\
Client & Admin console: the policy \code{WebAuthentication}\allowbreak\code{RemoteDesktop}\allowbreak\code{AllowedOrigins} lists the CRD client, \code{https://}\allowbreak\code{remotedesktop.google.com} & Policy verified; value assumed \\
Host & Google Chrome, not Chromium & Verified \\
Host & Force-install the CRD Security Key Extension (\code{ExtensionInstallForcelist}) & Verified \\
Host & CRD host: nothing to set; on by default & Verified \\
Google & The public CRD client speaks \code{remote-webauthn} & Unknown \\
\bottomrule
\end{tabular}
\end{table}

\begin{pitfall}[The open question]
We could not verify that Google's public CRD client implements the
\code{remote-webauthn} channel. Its code is not public. The extension's store
text names only Google's internal client, and its manifest lets only
internal pages message it~\cite{crdext} (verified). An administrator can test
it in a few minutes: set the policies of Table~\ref{tab:crd}, open
\code{webauthn.io} in the CRD session, and authenticate. If the prompt
appears on the Chromebook, the path works.
\end{pitfall}

\subsection{Compared with \code{ffwd-usbfs}}

\begin{itemize}
\item \textbf{Simpler and safer to run.} No Linux container, no
      \code{sudo}, no setuid sink, and no \code{uhid}. The key never leaves
      ChromeOS.
\item \textbf{WebAuthn only.} It serves web sites in the host's Chrome.
      SSH keys of type \code{sk}, OpenPGP and PIV on the host do not use
      Chrome, so they are not covered. Approach D still serves those.
\item \textbf{Same trust problem.} The client must trust the origin that the
      host reports, so a compromised host can ask for an assertion for any
      site. The specification says so~\cite{webauthn-ed}, and that is why
      Chrome allows it only by policy. The relay has the same weakness
      (Section~\ref{sec:security}).
\item \textbf{Depends on Google.} It needs a managed account, Google Chrome,
      and Google's closed CRD client.
\end{itemize}

\subsection{The older security key channel}

CRD has a second, older channel, \code{security-key}. It carries SSH agent
messages, not WebAuthn, and gives the session a socket in
\code{SSH\_AUTH\_SOCK}~\cite{crdgnubby} (verified). The policy
\code{RemoteAccessHostAllowGnubbyAuth} controls it; the policy is
internal to Google, and the client side is not public~\cite{crdpol}.
